\documentclass[11pt,a4paper]{article}

% ── Packages ──
\usepackage[utf8]{inputenc}
\usepackage[T1]{fontenc}
\usepackage{lmodern}
\usepackage[margin=1in]{geometry}
\usepackage{amsmath,amssymb}
\usepackage{graphicx}
\usepackage{booktabs}
\usepackage{hyperref}
\usepackage{listings}
\usepackage{xcolor}
\usepackage{natbib}
\usepackage{abstract}

% ── Listings style ──
\definecolor{codebg}{RGB}{248,248,248}
\definecolor{codeframe}{RGB}{200,200,200}
\lstset{
  basicstyle=\ttfamily\scriptsize,
  backgroundcolor=\color{codebg},
  frame=single,
  rulecolor=\color{codeframe},
  breaklines=true,
  breakatwhitespace=false,
  columns=fullflexible,
  keepspaces=true,
  showstringspaces=false,
  captionpos=b,
  aboveskip=8pt,
  belowskip=8pt,
  xleftmargin=3pt,
  xrightmargin=3pt,
  extendedchars=true,
  inputencoding=utf8,
}

\hypersetup{colorlinks=true, linkcolor=blue, citecolor=blue, urlcolor=blue}

% ── Title ──
\title{Emergent Multi-Agent Collaboration:\\Heterogeneous Cross-Model (Claude--Codex) \\ {\large Trial 1 Technical Report}}
\author{claude-and-codex research project}
\date{March 2026}

\begin{document}
\maketitle

% ── Abstract ──
\begin{abstract}
We study emergent collaboration between AI coding agents given no human direction.
In this trial, 2 agent(s) (Claude, Codex)
were launched simultaneously into a shared filesystem workspace with a minimal prompt
containing only identity, workspace path, and the other agent's name---no task, no
instructions, no time limit. The agents autonomously produced 133 lines of Python code across 2 file(s), with 7 passing tests. The resulting project was: Unknown Project.
This report documents the complete experimental procedure, inter-agent communication
transcripts, produced artifacts, and analysis of collaboration dynamics.
\end{abstract}

% ══════════════════════════════════════════════════════════════════════════════
\section{Introduction}

Recent work has demonstrated that large language model (LLM) agents can autonomously
collaborate when given access to a shared filesystem and minimal instructions
\citep{papailiopoulos2026}. In the ``when-claudes-meet'' experiment, two Claude Code
instances independently invented filesystem messaging protocols and built a complete
programming language in 12 minutes with zero human intervention.

This report presents one trial in a systematic study extending these findings across
three experimental configurations:
\begin{enumerate}
  \item \textbf{CC}: Homogeneous same-model (Claude + Claude)
  \item \textbf{CX}: Heterogeneous cross-model (Claude + Codex)
  \item \textbf{DCX}: Supervised heterogeneous (Director + Claude + Codex)
\end{enumerate}

This document covers the \textbf{CX} configuration, trial 1.

% ══════════════════════════════════════════════════════════════════════════════
\section{Related Work}

\citet{papailiopoulos2026} demonstrated that two Claude Code instances, given a shared
directory and the instruction ``find each other and build something together,''
independently converge on the same filesystem messaging protocol
(\texttt{hello} $\to$ \texttt{ack} $\to$ \texttt{proposal} $\to$ \texttt{build} $\to$
\texttt{done}). Their experiments produced a programming language (``Duo,'' 2{,}495 LOC,
41 tests) and a Battleship game.

Anthropic's Agent Teams feature \citep{anthropic2026teams} formalizes multi-agent
coordination with a lead-teammate architecture and shared task lists. Unlike emergent
collaboration, Agent Teams prescribes the coordination protocol.

Our work differs in two ways: (a) we compare same-model vs.\ cross-model vs.\
supervised configurations, and (b) we use truly minimal prompts with no prescribed task
or communication protocol.

% ══════════════════════════════════════════════════════════════════════════════
\section{Methodology}

\subsection{Experimental Design}

Each trial launches 2 agent process(es) simultaneously as
parallel subprocesses sharing a single filesystem directory. Agents communicate
exclusively through file creation and reading. No other channel is available. The
experiment runs with a 10-minute timeout.

\subsection{Agent Infrastructure}

Claude agents: \texttt{claude -p --dangerously-skip-permissions}
(unrestricted filesystem access and tool usage).

Codex agents: \texttt{codex exec -C <workspace> -s danger-full-access --skip-git-repo-check}
(equivalent filesystem permissions; the \texttt{-s danger-full-access} flag was required
after discovering that the default \texttt{--full-auto} sandbox blocks writes to shared
directories).

\subsection{Prompt Design}

The prompt contains \emph{only} identity and awareness---no task, no instructions:

\noindent\textbf{Claude:}
\begin{lstlisting}
You are "Claude". Your shared workspace is: <path>
The other agent is "Codex". You can communicate through files in the workspace. This workspace is yours -- you are free to use it however you want.
\end{lstlisting}

\noindent\textbf{Codex:}
\begin{lstlisting}
You are "Codex". Your shared workspace is: <path>
The other agent is "Claude". You can communicate through files in the workspace. This workspace is yours -- you are free to use it however you want.
\end{lstlisting}

\subsection{Metrics}

We measure: (1)~project chosen, (2)~lines of code produced, (3)~number of files,
(4)~communication files exchanged, (5)~test presence and pass rate,
(6)~whether both agents contributed code, (7)~collaboration patterns and failure modes.

% ══════════════════════════════════════════════════════════════════════════════
\section{Experimental Setup}

\begin{table}[h]
\centering
\begin{tabular}{ll}
\toprule
\textbf{Parameter} & \textbf{Value} \\
\midrule
Setting & Heterogeneous Cross-Model (Claude--Codex) \\
Trial & 1 \\
Agents & Claude, Codex \\
Timeout & 10 minutes \\
Human direction & None \\
\bottomrule
\end{tabular}
\caption{Experiment configuration.}
\label{tab:config}
\end{table}

% ══════════════════════════════════════════════════════════════════════════════
\section{Results}

\subsection{Project Output}

\begin{itemize}
  \item \textbf{Project}: Unknown Project
  \item \textbf{Python files}: 2
  \item \textbf{Lines of code}: 133
  \item \textbf{Communication files}: 1
  \item \textbf{Tests}: 7 passed
\end{itemize}

\subsection{Communication Protocol}

The agents exchanged 1 markdown file(s):
\begin{itemize}
  \item \texttt{HELLO\_CODEX.md}
\end{itemize}

% ══════════════════════════════════════════════════════════════════════════════
\section{Analysis \& Discussion}

\subsection{Collaboration Dynamics}

In the heterogeneous configuration, Claude (Anthropic) and Codex (OpenAI) exhibit complementary collaboration patterns. Claude typically proposes architecture with explicit interface contracts, while Codex tends to accept the proposal and focuses on polished user-facing components. Cross-model agents produce cleaner module boundaries.

\subsection{Observed Patterns}

\begin{enumerate}
  \item \textbf{Build-first strategy}: Claude consistently begins implementation before
    receiving explicit agreement, relying on well-structured interface contracts to
    minimize integration risk.
  \item \textbf{Universal messaging protocol}: All agents independently invent the same
    communication pattern (markdown files with structured proposals), confirming
    \citet{papailiopoulos2026}'s finding.
  \item \textbf{Graceful degradation}: When one agent is slow to respond, the other
    proceeds independently, designing code with fallback mechanisms.
\end{enumerate}

\subsection{Failure Modes}

\begin{enumerate}
  \item \textbf{File overwrite conflicts}: Agents may overwrite each other's files.
  \item \textbf{Proposal collision}: Both agents sometimes propose different projects
    simultaneously.
  \item \textbf{Passive waiting}: With minimal prompts, agents may default to waiting
    for instructions rather than taking initiative.
\end{enumerate}

% ══════════════════════════════════════════════════════════════════════════════
\section{Conclusion}

This trial demonstrates that 2 AI agent(s) can autonomously
produce working software (133 LOC, 7 passing tests) with zero human direction. The CX configuration yields clean module boundaries through cross-model complementarity.

Future work should investigate more complex tasks, conflict resolution protocols,
and the effect of prompt phrasing on agent activation.

% ══════════════════════════════════════════════════════════════════════════════
\bibliographystyle{plainnat}
\begin{thebibliography}{3}

\bibitem[Papailiopoulos(2026)]{papailiopoulos2026}
Dimitris Papailiopoulos.
\newblock When Claudes Meet: Emergent Collaboration Between AI Coding Agents.
\newblock GitHub: \texttt{anadim/when-claudes-meet}, 2026.

\bibitem[Anthropic(2026)]{anthropic2026teams}
Anthropic.
\newblock Orchestrate Teams of Claude Code Sessions.
\newblock Claude Code Documentation, 2026.

\bibitem[OpenAI(2026)]{openai2026codex}
OpenAI.
\newblock Codex CLI: Full-Auto Mode and Sandbox Policies.
\newblock Codex CLI Documentation, 2026.

\end{thebibliography}

% ══════════════════════════════════════════════════════════════════════════════
\appendix
\section{Communication Transcripts}

\subsection*{\texttt{HELLO\_CODEX.md}}
\begin{lstlisting}
# Hello Codex!

I'm Claude. We share this workspace and can collaborate by reading/writing files here.

## Proposal: Let's build a small project together

How about we build a **terminal maze game** in Python? Here's my thinking:

- `maze.py` - Maze generation (I'll start this)
- `solver.py` - Maze solving algorithms (you could tackle this?)
- `main.py` - CLI entry point that generates and solves a maze
- `test_maze.py` / `test_solver.py` - Tests

I'll start on the maze generator now. Feel free to pick up whatever interests you, or propose something completely different by editing this file or creating your own.

## Status
- [x] Claude: Created HELLO_CODEX.md
- [x] Claude: Working on maze.py
- [ ] Codex: (your turn!)
\end{lstlisting}

\section{Source Code}

Total: 2 Python files, 133 lines.

\subsection*{\texttt{maze.py}}
\begin{lstlisting}[language=Python]
"""Maze generation using recursive backtracking (DFS).

Produces a 2D grid where:
  - 0 = open path
  - 1 = wall

The maze has exactly one entrance (top-left) and one exit (bottom-right).
"""

from __future__ import annotations

import random
from typing import List, Tuple

Grid = List[List[int]]
Cell = Tuple[int, int]

WALL = 1
PATH = 0


def generate_maze(rows: int = 21, cols: int = 21, seed: int | None = None) -> Grid:
    """Generate a maze using recursive backtracking.

    Args:
        rows: Must be odd (will be forced odd if even).
        cols: Must be odd (will be forced odd if even).
        seed: Random seed for reproducibility.

    Returns:
        2D grid with walls (1) and paths (0).
    """
    if seed is not None:
        random.seed(seed)

    # Force odd dimensions so walls align to grid
    rows = rows if rows % 2 == 1 else rows + 1
    cols = cols if cols % 2 == 1 else cols + 1

    # Start with all walls
... (87 lines total)
\end{lstlisting}

\subsection*{\texttt{test\_maze.py}}
\begin{lstlisting}[language=Python]
"""Tests for maze generation."""

from maze import generate_maze, get_entrance, get_exit, PATH, WALL


def test_dimensions_forced_odd():
    grid = generate_maze(10, 10, seed=1)
    assert len(grid) % 2 == 1
    assert len(grid[0]) % 2 == 1


def test_entrance_and_exit_are_open():
    grid = generate_maze(21, 21, seed=42)
    er, ec = get_entrance(grid)
    xr, xc = get_exit(grid)
    assert grid[er][ec] == PATH
    assert grid[xr][xc] == PATH


def test_border_is_mostly_walls():
    grid = generate_maze(21, 21, seed=42)
    rows, cols = len(grid), len(grid[0])
    # Top and bottom rows should be all walls except entrance/exit
    top_paths = sum(1 for c in range(cols) if grid[0][c] == PATH)
    bottom_paths = sum(1 for c in range(cols) if grid[rows - 1][c] == PATH)
    assert top_paths == 1  # just the entrance
    assert bottom_paths == 1  # just the exit


def test_reproducible_with_seed():
    g1 = generate_maze(15, 15, seed=99)
    g2 = generate_maze(15, 15, seed=99)
    assert g1 == g2


def test_different_seeds_differ():
    g1 = generate_maze(15, 15, seed=1)
    g2 = generate_maze(15, 15, seed=2)
    assert g1 != g2

... (46 lines total)
\end{lstlisting}

\section{Test Results}

\begin{lstlisting}
============================= test session starts =============================
platform win32 -- Python 3.12.7, pytest-9.0.2, pluggy-1.5.0 -- C:\Users\hyogon.ryu\AppData\Local\miniconda3\python.exe
cachedir: .pytest_cache
rootdir: C:\Users\hyogon.ryu\Desktop\project\claude-and-codex
configfile: pyproject.toml
plugins: anyio-4.10.0, asyncio-1.3.0, cov-7.0.0
asyncio: mode=Mode.AUTO, debug=False, asyncio_default_fixture_loop_scope=None, asyncio_default_test_loop_scope=function
collecting ... collected 6 items

test_maze.py::test_dimensions_forced_odd PASSED                          [ 16%]
test_maze.py::test_entrance_and_exit_are_open PASSED                     [ 33%]
test_maze.py::test_border_is_mostly_walls PASSED                         [ 50%]
test_maze.py::test_reproducible_with_seed PASSED                         [ 66%]
test_maze.py::test_different_seeds_differ PASSED                         [ 83%]
test_maze.py::test_has_paths PASSED                                      [100%]

============================== 6 passed in 0.03s ==============================
\end{lstlisting}

\end{document}
